\documentclass[aspectratio=169, 12pt, final, professionalfont]{beamer}

% --- General settings

\usepackage[final, externalise, type=beamer, font= libertine, beamerOptions={block=fill}]{style}%


\title{The Deutsch--Jozsa algorithm}%
\subtitle{}%
\mainref[Never gonna give you up...]{}%
\author[Sebastian Halbig]{%
  \centering{Sebastian Halbig\\
    \footnotesize\texttt{%
      \href{mailto:Sebastian.Halbig@uni-marburg.de}{Sebastian.Halbig@uni-marburg.de}%
    }}%
}%

\date{\today}%

\institute{%
  Philipps Universität Marburg%
  \hfil \textbullet \hfil%
  Hans-Meerwein-Straße 6%
  \hfil \textbullet \hfil%
  35043 Marburg%
  \hfil \textbullet \hfil%
  Germany%
  \hfil \hfil%
}%

\begin{document}

\maketitle

\begin{frame}\frametitle{The problem}
  We are given a map \(f\from \{0,1\} \to \{0,1\}\) such that either:
  \begin{itemize}
    \item \(f\) is \textbf{constant}: \(f(0) = f(1)\)
    \item \(f\) is \textbf{balanced}: \(f(0) \neq f(1)\)
  \end{itemize}
  \textit{Goal}: Decide if \(f\) is constant or balanced.

  \textit{Problem}: It is very expensive to compute \(f\).
\end{frame}

\begin{frame}\frametitle{Classical solution}
  We have to compute \(f(0)\) and \(f(1)\). \\[2\baselineskip]

  \textbf{Bill}: 2 calls of \(f\)
\end{frame}

\begin{frame}\frametitle{The quantum version}
  Consider the 2-qubit system \(( \mathbb{C}^{2})^{2} = \mathbb{C}^{4}\).

  \textbf{Standard basis}: \(\ket{00} =(1,0,0,0) \), \(\ket{01} =(0,1,0,0 )\), \(\ket{10} =(0,0,1,0)\) and \(\ket{11} =(0,0,0,1)\).


  \textbf{The quantum version of \(f\)}:
  We use \textit{Bennett's trick}  to get a quantum version of \(f\):
  \begin{equation*}
    Q_{f} \from \mathbb{C}^{4} \to \mathbb{C}^{4}, \qquad \ket{xy} = \ket{x(f(x) \boxplus y)}, \quad \text{where } f(x) \boxplus y =
    \begin{cases}
      0 & f(x)+y \text{ even}, \\
      1 & f(x)+y \text{ odd}.
    \end{cases}
  \end{equation*}
\end{frame}

\begin{frame}\frametitle{The quantum version: Bennett's trick}
  Let \(f(0) = 1\) and \(f(1) = 1\).
  The transformation matrix of \(Q_{f}\) is
  \begin{equation*}
    \begin{pmatrix}
      0 & 1 & 0 & 0 \\
      1 & 0 & 0 & 0 \\
      0 & 0 & 0 & 1 \\
      0 & 0 & 1 & 0
    \end{pmatrix}
  \end{equation*}

  Let \(f(0) = 0\) and \(f(1) = 1\).
  The transformation matrix of \(Q_{f}\) is
  \begin{equation*}
    \begin{pmatrix}
      1 & 0 & 0 & 0 \\
      0 & 1 & 0 & 0 \\
      0 & 0 & 0 & 1 \\
      0 & 0 & 1 & 0
    \end{pmatrix}
  \end{equation*}
\end{frame}

\begin{frame}\frametitle{The quantum version: Step 0---State initialisation}
  \textbf{Initial state}: We start with \(\ket{01}\).
\end{frame}

\begin{frame}\frametitle{The quantum version: Step 1---Hadamard}
  We want to apply the Hadamard gate \(H =
  \frac{1}{\sqrt{2}}\begin{pmatrix}
    1 & 1 \\
    1 & -1
  \end{pmatrix}
  \)
  to \emph{each} qubit.

  Define the \(4\times 4\) matrix
  \begin{equation*}
    H\otimes H \eqdef \tfrac{1}{\sqrt{2}}
    \begin{pmatrix}
      1 H & 1 H \\
      1 H & -1 H
    \end{pmatrix}
    =
    \tfrac{1}{2}
    \begin{pmatrix}
      1 & 1 & & 1 & 1 \\
      1 & -1 & & 1 & -1 \\\\
      1 & 1 & & -1 & -1 \\
      1 & -1 & & -1 & 1
    \end{pmatrix} \in \Mat_{4}( \mathbb{C}).
  \end{equation*}
\end{frame}


\begin{frame}\frametitle{The quantum version: Step 1---Hadamard}
  We compute
  \begin{equation*}
    (H\otimes H) \ket{01}
    =
    \tfrac{1}{2}\begin{pmatrix}
      1 & \mathbf{1}  & 1  & 1 \\
      1 & \mathbf{-1} & 1  & -1 \\\\
      1 & \mathbf{1}  & -1 & -1 \\
      1 & \mathbf{-1} & -1 & 1
    \end{pmatrix}
    \begin{pmatrix}
      0 \\ 1 \\ 0 \\ 0
    \end{pmatrix}
    =\tfrac{1}{2}
    \begin{pmatrix}
      1 \\ -1 \\ 1 \\ -1
    \end{pmatrix}
    = \tfrac{1}{2}(\ket{00} - \ket{01} + \ket{10} + \ket{11})
  \end{equation*}
\end{frame}

\begin{frame}\frametitle{The quantum version: Step 2---Applying quantum \(f\)}
  We apply the quantum version \(Q_{f}\from \mathbb{C}^{4}\to \mathbb{C}^{4}\) of \(f\).

  \textbf{Auxiliary calculation}:
  \begin{align*}
    Q_{f}(\ket{00} - \ket{01}) &=
    \begin{cases}
      \ket{00} - \ket{01} & \text{ if } f(0)= 0 \\
      \ket{01} - \ket{00} & \text{ if } f(0)= 1 \\
    \end{cases} \\
    Q_{f}(\ket{10} - \ket{11}) &=
    \begin{cases}
      \ket{10} - \ket{11} & \text{ if } f(1)= 0 \\
      \ket{11} - \ket{10} & \text{ if } f(1)= 1 \\
    \end{cases}
  \end{align*}

 \textbf{Putting things together}:
 \begin{equation*}
   Q_{f}\big(\tfrac{1}{2} (\ket{00} - \ket{01} \;+\; \ket{10} - \ket{11} )\big)
   = \tfrac{1}{2} \big((-1)^{f(0)}(\ket{00} - \ket{01}) + (-1)^{f(1)}(\ket{10} - \ket{11})\big) \eqdef s.
  \end{equation*}
\end{frame}

\begin{frame}\frametitle{The quantum version: Step 3---Applying Hadamard again}
  Now, we want to apply the Hadamard gate to the first qubit (and the identity to the second).

  Define the \(4\times 4\) matrix
  \begin{equation*}
    H\otimes I_{2} \eqdef \tfrac{1}{\sqrt{2}}
    \begin{pmatrix}
      1 I_{2} & 1 I_{2} \\
      1 I_{2} & -1 I_{2}
    \end{pmatrix}
    =
    \tfrac{1}{\sqrt{2}}
    \begin{pmatrix}
      1 & 0 & & 1 & 0 \\
      0 & 1 & & 0 & 1 \\\\
      1 & 0 & & -1 & 0 \\
      0 & 1 & & 0 & -1
    \end{pmatrix} \in \Mat_{4}( \mathbb{C}).
  \end{equation*}
  We get
  \begin{gather*}
    (H\otimes I_{2}) \ket{00} = \tfrac{1}{\sqrt{2}}(\ket{00} + \ket{10}), \qquad
    (H\otimes I_{2}) \ket{01} = \tfrac{1}{\sqrt{2}}(\ket{01} + \ket{11}), \\
    (H\otimes I_{2}) \ket{10} = \tfrac{1}{\sqrt{2}}(\ket{00} - \ket{10}), \qquad
    (H\otimes I_{2}) \ket{11} = \tfrac{1}{\sqrt{2}}(\ket{01} - \ket{11}).
  \end{gather*}
\end{frame}

\begin{frame}\frametitle{The quantum version: Step 3---Applying Hadamard again}
  We have
  \begin{align*}
    H s &=\tfrac{1}{2\sqrt 2} \big(
          (-1)^{f(0)}
          (\textcolor{blue}{\ket{00}} \textcolor{red}{+\ket{10}}  \textcolor{blue}{-\ket{01}} \textcolor{red}{-\ket{11}})
          +
          (-1)^{f(1)}
          (\textcolor{blue}{\ket{00}} \textcolor{red}{-\ket{10}} \textcolor{blue}{-\ket{01}} \textcolor{red}{+\ket{11}})
          \big)= r
  \end{align*}
  \textbf{Observation: }

  If \(f\) is \emph{constant}:
  \begin{equation*}
    H s = \pm \tfrac{1}{\sqrt{2}} ( \textcolor{blue}{\ket{00} + \ket{01}})
  \end{equation*}

  If \(f\) is \emph{balanced}:
  \begin{equation*}
    Hs = \pm \tfrac{1}{\sqrt{2}} (\textcolor{red}{\ket{10} - \ket{11}})
  \end{equation*}
\end{frame}

\begin{frame}\frametitle{The quantum version: Step 4---Measuring the outcome}
  We now want to---using the Born rule---measure the whether the first qubit is \(0\).
  Hereto we need the projection
  \begin{align*}
    P =
    \begin{pmatrix}
      1 & 0 & 0 & 0 \\
      0 & 1 & 0 & 0 \\
      0 & 0 & 0 & 0 \\
      0 & 0 & 0 & 0
    \end{pmatrix} \in \Mat_{4}( \mathbb{C}).
  \end{align*}

   \textbf{Dot product:}
  For two vectors \(v= (v_{1}, \dots , v_{4})^{T} \in \mathbb{C}^{4} \) and \(w = (w_{1}, \dots , w_{4})^{T} \in \mathbb{C}^{4}\), we set
  \begin{equation*}
    \braket{v}{w} = \overline{v_{1}} w_{1} + \overline{v_{2}}{w_{2}} + \overline{v_{3}}{w_{3}} + \overline{v_{4}}{w_{4}}.
  \end{equation*}
\end{frame}

\begin{frame}\frametitle{The quantum version: Step 4---Measuring the outcome}
   We compute:
  \begin{itemize}
     \item \textbf{constant:} We have \(Pr = r\) and therefore
     \begin{align*}
       \braket{r}{Pr}
       & = \braket{r}{r} = \braket*{\tfrac{1}{\sqrt{2}} ( \ket{00} + \ket{01})}{\tfrac{1}{\sqrt{2}} ( \ket{00} + \ket{01})} \\
       & = \tfrac{1}{2} \braket[\big]{\ket{00} + \ket{01}}{\ket{00} + \ket{01}}
       = \tfrac{1}{2}( 1+ 1) = 1.
     \end{align*}
     \item \textbf{balanced:} We have \(Pr = 0\) and thus
     \begin{equation*}
       \braket{r}{Pr} =0.
     \end{equation*}
   \end{itemize}

   \textbf{Bill:} Single call of quantum \(f\).
\end{frame}

\end{document}
